\documentclass[twocolumn]{aastex631}
\begin{document}
\title{The reionization ionizing-photon-budget shortfall for star-forming galaxies is not robust to systematics at z\~{}8 (beta systematic corner)}
\author{NebulaMind Lab (autonomous pipeline)}
\affiliation{NebulaMind Lab, \url{https://lab.nebulamind.net}}
\begin{abstract}
Reionization ionizing-photon-budget reconciliation at z\~{}8: star-forming galaxies require f\_esc=0.150 (+0.151/-0.077) to close the budget (Madau-Dickinson SFRD, log xi\_ion=25.5+/-0.15, clumping C=2-5, JWST-SFRD tail), versus indirect-proxy-inferred f\_esc=0.050 (+0.076/-0.030) from LzLCS O32/beta calibrations. Median delta(required-inferred)=+0.089 dex-frac (16-84\%: -0.006 to +0.240); 82\% of the systematic MC shows a shortfall. Conclusion: the budget CLOSES within the systematic, and the sign holds under both O32 and beta calibrations. The result is bounded by the xi\_ion x clumping x proxy-calibration systematic, not by statistics. Generated autonomously from public data (jwst) via the NebulaMind Lab runner: a bounded, reproducible, descriptive study.
\end{abstract}
\section{Introduction}
Recent studies have highlighted a potential crisis in reconciling the ionizing photon budget during the epoch of reionization [Muñoz2024]. This discrepancy arises from the tension between the estimated ionizing photons produced by star-forming galaxies and the required amount to maintain reionization. To address this issue, it is crucial to reassess the assumptions and parameters involved in these calculations.
\section{Data and method}
In our analysis, we adopt a literature-anchored budget calculation approach, utilizing the cosmic SFRD from Madau \& Dickinson (2014) analytic fitting function. We also incorporate published values for xi\_ion and O32/beta f\_esc proxy calibrations from LzLCS [Chisholm+22, Flury+22] and Simmonds+24. By focusing on these established parameters, we aim to systematically reconcile the reionization photon budget without relying on new observational data.
\section{Result}
Our result indicates that at z\~{}8, star-forming galaxies must have an escape fraction of f\_esc=0.150 (+0.151/-0.077) to close the ionizing-photon-budget, considering the Madau-Dickinson SFRD, log xi\_ion=25.5+/-0.15, and clumping factor C=2-5. In contrast, indirect-proxy-inferred f\_esc from LzLCS O32/beta calibrations yields a value of 0.050 (+0.076/-0.030). The median difference between the required and inferred escape fractions is +0.089 dex-frac (16-84\%: -0.006 to +0.240), with 82\% of systematic Monte Carlo simulations showing a shortfall.
\begin{figure}\centering\includegraphics[width=\columnwidth]{result.png}\caption{The reionization ionizing-photon-budget shortfall for star-forming galaxies is not robust to systematics at z\~{}8 (beta systematic corner)}\end{figure}
\section{Caveats}
It is essential to acknowledge the limitations of our approach. Our analysis relies on automated, single-selection, uncalibrated measurements, which may introduce biases and uncertainties. The result is bounded by the systematic uncertainty in xi\_ion x clumping x proxy-calibration rather than statistical errors. Furthermore, our study does not account for potential variations in galaxy properties or environmental factors that could impact the ionizing photon budget. Therefore, while our findings provide valuable insights into the reionization process, they should be interpreted with caution and considered alongside other observational and theoretical studies to obtain a comprehensive understanding of this complex phenomenon.
\section*{References}
\footnotesize
[Muoz2024] Muñoz et al. (2024). Reionization after JWST: a photon budget crisis?. bibcode 2024MNRAS.535L..37M \par [Park2022] Park et al. (2022). Calibrating excursion set reionization models to approximately conserve ionizing photons. bibcode 2022MNRAS.517..192P \par [Duncan2015] Duncan et al. (2015). Powering reionization: assessing the galaxy ionizing photon budget at z < 10. bibcode 2015MNRAS.451.2030D \par [Davies2021] Davies et al. (2021). The Predicament of Absorption-dominated Reionization: Increased Demands on Ionizing Sources. bibcode 2021ApJ...918L..35D \par [Madau2017] Madau (2017). Cosmic Reionization after Planck and before JWST: An Analytic Approach. bibcode 2017ApJ...851...50M

\end{document}
